\documentclass[12pt]{article}
\usepackage{amsmath,amssymb,amsfonts}
\usepackage[margin=1in]{geometry}

\begin{document}

\textbf{Chapter 8, Problem 3.} Minkowski's Inequality.

\textbf{(a)} Show that $|f(x)+g(x)|^p \le |f(x)+g(x)|^{p-1}|f(x)| + |f(x)+g(x)|^{p-1}|g(x)|$, and hence that if $f$ and $g$ belong to $L_p$ (the space of $p$-integrable functions), then so does $f+g$. Thus $L_p$ is indeed a vector space.

\textbf{(b)} Show that for $p > 1$,
\[
\int_a^b |f+g|^p\,dx \le \int_a^b |f+g|^{p-1}|f|\,dx + \int_a^b |f+g|^{p-1}|g|\,dx,
\]
and that $|f+g|^{p-1}$ belongs to $L_q$ if $f$ and $g$ belong to $L_p$, where $1/p+1/q=1$.

\textbf{(c)} Apply H\"older's inequality (see Problem 8.2) to each of the two integrals on the right-hand side of the inequality in part (b), and hence show that
\[
\left[\int_a^b|f+g|^p\,dx\right]^{1/p} \le \left[\int_a^b|f|^p\,dx\right]^{1/p}+\left[\int_a^b|g|^p\,dx\right]^{1/p}.
\]
This is Minkowski's inequality.

\bigskip
\textbf{Solution.}

\textbf{Part (a):}
By the triangle inequality,
\[
|f(x)+g(x)| \le |f(x)| + |g(x)|.
\]
Multiplying both sides by $|f(x)+g(x)|^{p-1}\ge 0$:
\[
|f+g|^p \le |f+g|^{p-1}|f| + |f+g|^{p-1}|g|.
\]
Now if $f,g\in L_p$, then $|f|^p$ and $|g|^p$ are integrable. Since $|f+g|\le|f|+|g|$:
\[
|f+g|^p \le (|f|+|g|)^p \le 2^p\max(|f|^p,|g|^p) \le 2^p(|f|^p+|g|^p).
\]
Integrating shows $\int|f+g|^p\,dx < \infty$, so $f+g\in L_p$. Hence $L_p$ is closed under addition and is a vector space. $\checkmark$

\textbf{Part (b):}
Integrating the inequality from part (a):
\[
\int_a^b |f+g|^p\,dx \le \int_a^b |f+g|^{p-1}|f|\,dx + \int_a^b |f+g|^{p-1}|g|\,dx.
\]
Now, since $1/p+1/q=1$ implies $q = p/(p-1)$, we have $(p-1)q = p$. Thus
\[
\int_a^b \left(|f+g|^{p-1}\right)^q dx = \int_a^b |f+g|^p\,dx < \infty,
\]
since $f+g\in L_p$ by part (a). Therefore $|f+g|^{p-1}\in L_q$. $\checkmark$

\textbf{Part (c):}
Apply H\"older's inequality to each integral on the right:
\[
\int |f+g|^{p-1}|f|\,dx \le \left[\int|f|^p\,dx\right]^{1/p}\left[\int|f+g|^{(p-1)q}\,dx\right]^{1/q}
= \|f\|_p\left[\int|f+g|^p\,dx\right]^{1/q},
\]
and similarly
\[
\int|f+g|^{p-1}|g|\,dx \le \|g\|_p\left[\int|f+g|^p\,dx\right]^{1/q}.
\]
Combining:
\[
\int|f+g|^p\,dx \le (\|f\|_p+\|g\|_p)\left[\int|f+g|^p\,dx\right]^{1/q}.
\]
Dividing both sides by $\left[\int|f+g|^p\,dx\right]^{1/q}$ (assuming it is nonzero; if zero, the result is trivial):
\[
\left[\int|f+g|^p\,dx\right]^{1-1/q} \le \|f\|_p+\|g\|_p.
\]
Since $1-1/q = 1/p$:
\[
\boxed{\|f+g\|_p \le \|f\|_p + \|g\|_p.}
\]
This is Minkowski's inequality, confirming that $\|\cdot\|_p$ defines a norm on $L_p$. $\blacksquare$

\end{document}
